\section{Calibration of the primary coefficient transition}
\label{sec:primary-calibration-producer}
The primary transition must do more than bound a regular insertion. Its
coefficient must approach the prescribed linear transport on the finite
coefficient frame. We prove this by the same three comparisons used in
the construction: the actual principal term to its Gaussian reference,
the Gaussian reference to the linear predicted field, and that field to
the finite polynomial probe calculation.

Fix a frame of total degree $d$ and a nonconstant coordinate $i$. Write
$g_i\in\{1,\ldots,d\}$ for its spatial grade and $\ell_i$ for its
order of vanishing in the field amplitude. These are different gradings.
The primary insertion $h_i$ is the local section at fine threshold
$\theta_f$, normalized by $\theta_f^{-\ell_i}$ as in
Section~\ref{sec:local-sections}. Its regularity price is bounded by the
fixed fine-frame constant $T_f$. All constants in this section may depend
on the chosen frame and the block factor $L$. Periods and the incoming
coupling remain uniform. Both adjacent kernel levels are at least four,
and the fine level is at least $d$.

\subsection{Removing the Gaussian fluctuation at the flat field}
Let $h$ be one real component of $h_i$. Denote by $H(V)$ its reference
Gaussian integral and by
\[
 h_0(V)=h\left(\exp\bigl(\mathcal A\,
                  \operatorname{crop}_{R}(\log V)\bigr)\right)
\]
the linear prediction, with the installed response $\mathcal A$ and
$R=45(1+B_c)$. This formula is used as a germ at $V=1$; its crop is in
the identity gauge. Let $k$ be the derivative order of a frame probe.
Purchase regularity through order $n\ge k+1$, as well as $n\ge2$.
The additional derivative will control the change of the insertion's
base point.

Use the fixed Gaussian measure at $V=1$ and let $R_V(z)$ be the
normalized density ratio. Put $A_V(z)=h(W(V,z))$, where $W$ is the
actual received reference configuration. With the cutoff $\Gamma$
of the preceding section, the exact difference of integrands is
\begin{equation}\label{eq:primary-integrand-difference}
 \Gamma R_V A_V-h_0(V)
 =\Gamma R_V(A_V-h_0(V))
       +\Gamma(R_V-1)h_0(V)+(\Gamma-1)h_0(V).
\end{equation}
Thus no derivative of the varying density is discarded.

Write $P(z)\ge\max\{1,\|\mathrm{tape}(z)\|\}$ for the fixed
Gaussian coordinate norm augmented by one. On the cutoff support,
the display comparison of Section~\ref{sec:principal-reference-attenuation}
gives, for every jet through order $n$, an error of fine direction price
\[
 p(z)=C_n\frac{Q s^2+g\|\mathrm{tape}(z)\|}{\theta_f}.
\]
The linear predictor's argument jets have bounded fine direction price
$W\ge1$. The nonlinear jets have price at most $W+p(z)$.
For a partition with $r$ arguments, subtract the two chain-rule terms
by first changing their base point and then their arguments one at a time.
Changing the base point costs at most
$C\tau p(z)(W+p(z))^r$: integrate the $(r+1)$st derivative of the
insertion along the healthy radial segment proved earlier. Changing the
arguments costs at most
$Cr\tau p(z)(W+p(z))^{r-1}$. All these terms use $r+1\le n$.
Summing the finitely many partitions proves
\[
 |D^j(A_V(z)-h_0(V))|_{V=1}
       \le C_n\tau p(z)(W+p(z))^n,
                         \qquad j\le k.
\]
The same argument without subtraction gives $|D^jh_0(1)|\le C_n\tau$.

The density-ratio calculation in
Section~\ref{sec:gaussian-differentiation} gives a small factor
$\theta_f$ times a polynomial in the read radius and $P(z)^2$
in every positive-order jet of $R_V-1$ at $1$.
Its value there is zero since $R_1=1$. More explicitly, differentiate
$\tfrac12\log\det C(V)-\tfrac12\langle z,C(V)z\rangle$
with the reference precision $C$. Each positive-order term contains
at least one admissible coarse direction, and the inverse and precision
jets bound it by $K_L\theta_f(1+B_c)^aP(z)^2$. The partition formula
for its exponential multiplies such factors, retaining at least one.

The cutoff error is controlled even outside its support. It is one on
the plateau $g\|\mathrm{tape}(z)_a\|\le\varepsilon\theta_f/8$
for every coordinate $a$. If a coordinate fails this test, then
$(g/\theta_f)P(z)^2\ge\varepsilon/8$, while $|\Gamma-1|\le2$.
Consequently
\[
 |\Gamma(z)-1|\le \frac{16}{\varepsilon}
                           \frac g{\theta_f}P(z)^2.
\]
This bound does not require evaluating $A_V$ off the cutoff support:
there the first two terms of \eqref{eq:primary-integrand-difference}
vanish identically.

Collect the displayed small factors into $J_\beta\ge0$ so that
\[
 J_\beta\le K_L(\theta_f+g/\theta_f)(1+B_c)^a,
 \qquad p(z)\le J_\beta P(z)^2.
\]
The bounded threshold ratio supplies $s^2/\theta_f\le
K_L\theta_f(1+B_c)^a$. Increase the coupling until $J_\beta\le1$.
Leibniz' rule in \eqref{eq:primary-integrand-difference} and the
partition bounds then give
\[
 |D^j(\Gamma R_V A_V-h_0(V))|_{V=1}
                \le C_n\tau J_\beta P(z)^{4(n+1)},
                         \qquad j\le k.
\]
All factors from higher-order products are absorbed by the displayed
power, since $P\ge1$ and $J_\beta\le1$. The fixed Gaussian moments
bound its integral by $K'_L(1+B_c)^b$. Compact cutoff support and the
common local smoothness justify differentiation under the integral;
the polynomial bound also dominates the subtracted constant term.
The monitor estimate from the preceding section implies
\begin{equation}\label{eq:primary-flat-gaussian}
 \max_{j\le k}|D^jH(1)-D^jh_0(1)|\longrightarrow0
                                      \quad(\beta\to\infty).
\end{equation}
The bound is uniform in every forest, replica choice and allowed
period. Averaging therefore preserves it. Combining it with the
actual-to-reference comparison gives the same convergence for the
actual principal insertion.

\subsection{Normalized stencil functions and their jets}
We give the finite-dimensional estimate needed to sample $h_0$ on the
frame probes. Let $H_i$ be the analytic unnormalized stencil function
of coordinate $i$, on the finite-dimensional space of its link logs.
It satisfies $D^jH_i(0)=0$ for $j<\ell_i$. For $0<t\le1$ set
\[
                       H_{i,t}(v)=t^{-\ell_i}H_i(tv).
\]
On any fixed ball $\|tv\|\le R_0$, Taylor's formula and compactness
give, for every prescribed order $j$,
\begin{equation}\label{eq:primary-normalized-jets}
                    \|D^jH_{i,t}(v)\|
                           \le C(1+\|v\|)^{\ell_i}.
\end{equation}
Indeed let $a=\max\{\ell_i-j,0\}$. The vanishing derivatives imply
$\|D^jH_i(tv)\|\le C\|tv\|^a$ by integrating $a$ times from zero.
For $a=0$, use the maximum of the continuous derivative on the ball.
The scaling factor is $t^{j-\ell_i+a}\le1$, with nonnegative exponent;
and $\|v\|^a\le(1+\|v\|)^{\ell_i}$. Taking a maximum over the
finitely many required orders proves the bound with one $C$.

Suppose two $C^k$ curves $v,q$ have jets at zero bounded by $U+e,U$,
respectively, and their jet differences through order $k$ are at most
$e\ge0$. Suppose also their values, multiplied by $t$, lie in the
fixed ball. Along the segment joining their values, use
\eqref{eq:primary-normalized-jets} through order $k+1$.
For a partition with $r\le j\le k$ blocks, the difference of the
outer $r$th derivatives costs at most
$C(1+U+e)^{\ell_i}e$. Replacing its $r$ arguments one at a time
costs at most $Cr(1+U+e)^{\ell_i+r}e$. Thus
\begin{equation}\label{eq:primary-stencil-stability}
 \max_{j\le k}\|D^j(H_{i,t}\circ v)(0)
                   -D^j(H_{i,t}\circ q)(0)\|
           \le C_{k,i,U}(1+e)^{\ell_i+k}e.
\end{equation}
This estimate is uniform as $t\downarrow0$. The same partition
calculation without subtraction bounds the jets of
$H_{i,t}\circ q$ uniformly. The constant is a finite sum over
partitions, not a limit at the singular normalization $t=0$.

\subsection{The cropped torus probe and the full polynomial forecast}
A probe consists of a fixed polynomial potential and a derivative
order $k\le d$. Its test-link curve is the holonomy curve defined in
Section~\ref{sec:local-sections}. At derivative order $j\le k$,
its link-log jet on the covering lattice is a polynomial field of
degree at most $dj$. This follows by expanding the exponential,
product and logarithm to order $j$: every term is a product of at
most $j$ potential factors, each of degree at most $d$. Integration
along a unit bond and finite differences do not increase that bound.
The Taylor coefficients are fixed by the probe.

Normalize the coarse probe by its threshold $y=\theta_c$ and the
fine link field by $t=\theta_f$, with $y/t\le2$ and $y\le1$.
Apply the installed response to the cropped torus link-log curve and
restrict to the fixed fine stencil; call the resulting normalized
curve $v_R$. Let $q$ be the full covering-lattice prediction obtained
by applying the same response to the polynomial link-log jets.
Only its $k$-jet is used; it is equivalently obtained by the finite
polynomial jet algebra of the probe. Both curves vanish at zero.
The exact correspondence between periodic and covering convolution
holds because the retained kernel is periodized from its absolutely
summable covering kernel.

On covering sites of distance at most $R$, the crop and the full
probe have identical jets. The room inequality $2R<c$ ensures that
these sites have their canonical lifts, so there is no seam error
there. Outside this region, the probe jets have polynomial growth
and the cropped jets have norm at most $K(1+4R)^{dk}$.
The response strip bound gives
\begin{equation}\label{eq:primary-probe-tail}
 \max_{j\le k}\|D^jv_R(0)-D^jq(0)\|
       \le C\{1+K(1+4R)^{dk}\}e^{-\sigma R/2}=:e_R.
\end{equation}
To verify the tail estimate, take half the exponential weight to
bound the polynomial factor uniformly, and sum the remaining
weighted kernel. The finite stencil offsets and four directions
supply only a fixed multiplier. The torus sum is the same sum
regrouped by residue class; absolute convergence permits this
regrouping. This is the cropped-versus-full comparison, including
all derivative orders, for the actual retained response.

The full predicted jets have a bound $U$ independent of the period
and $R$. Indeed convolution of a degree-$dj$ polynomial costs its
fixed weighted kernel moment. In a nonzero chain-rule term with
$r\ge1$ blocks, the threshold factor is $y^r/t\le y/t\le2$.
The finitely many probe partitions therefore have a common bound.
Apply \eqref{eq:primary-stencil-stability} with $t=\theta_f$ and
$R_0=1$: both input values are zero and hence lie in this fixed
positive-radius ball. Since $e_R\to0$, the resulting
normalized stencil jet difference tends to zero. Multiplication by
the fixed physical factor $L^d$ preserves that convergence.

To identify the prescribed coefficient transport explicitly, let
$\mathscr S_t$ be diagonal multiplication by $t^{\ell_i}$, let
$\mathscr D_L$ be diagonal multiplication by $L^{g_i}$, and let
$\mathscr V_L$ be diagonal multiplication by $L^{d-g_i}$. If
$\mathscr R$ is coefficient recovery of the full polynomial forecast,
the construction's heading is the matrix
\[
 \mathscr S_{\theta_c}\,\mathscr R\,
             \mathscr D_L\,\mathscr S_{\theta_f^{-1}},
 \qquad
 x_i=\mathscr S_{\theta_c}\mathscr R
       \mathscr D_L\mathscr S_{\theta_f^{-1}}\mathscr V_L e_i.
\]
On column $i$, the last three diagonal factors contribute
$L^{g_i}\theta_f^{-\ell_i}L^{d-g_i}
 =L^d\theta_f^{-\ell_i}$. The remaining coarse scaling substitutes
$\theta_c$ times the probe potential into its polynomial jet.
Consequently, for every frame probe $p$,
\[
 \operatorname{probe}_p(x_i)
           =L^d D^{k_p}(H_{i,\theta_f}\circ q_p)(0).
\]
This identity is an equality of the finite jet calculation: evaluate
the polynomial transport on the probe, substitute the response's
polynomial moments, and use the local section's exact coefficient
recovery. The normalization on each side is
$L^d\theta_f^{-\ell_i}$, with the same coarse threshold in the
probe. Thus \eqref{eq:primary-probe-tail} identifies the linear
prediction's probes with those of $x_i$, to an error tending to zero.
It does not replace the retained response by another continuum kernel.

\subsection{From probes to the paid primary coefficient}
Let $F_i$ be the actual complex principal insertion, already multiplied
by $L^d$, and let $b_i$ be its extracted coefficient vector. The
nonwrapping coarse stencil allows coefficient recovery by the fixed
finite probe family. Ordered derivative convergence in
\eqref{eq:primary-flat-gaussian} implies probe convergence: each
probe derivative is a finite combination of ordered derivatives
along its fixed test-link curve. The coefficients are uniformly
bounded by its finite local price. The last subsection and the
triangle inequality therefore give
\[
       \max_p|\operatorname{probe}_p(b_i-x_i)|\longrightarrow0.
\]
A finite separating family of linear probes controls the coefficient
norm: on its unit sphere their maximum has a positive minimum by
compactness. Equivalently use the finite recovery matrix from
Section~\ref{sec:local-sections}. Denote its bound by $T_{\rm test}$.
Then $\|b_i-x_i\|_1\le T_{\rm test}\delta_\beta$, where
$\delta_\beta\to0$ is the maximum displayed probe error.

The remaining regular rows have weighted sum tending to zero after
$L$ is fixed, by the nonprincipal estimate of
Section~\ref{sec:leaf-expansion}. The extraction map on them costs
the fixed read constant $T_{\rm read}$. If their summed majorant is
$u_\beta$, the total extracted coefficient differs from $x_i$ by at most
\[
                    p_\beta=T_{\rm test}\delta_\beta
                                  +T_{\rm read}u_\beta\longrightarrow0.
\]
The equality splitting the sum into its principal row and the other
rows is justified by their absolute weighted summability.

Move this coefficient error into the local-section remainder. The
local splitting identity of Section~\ref{sec:local-splitting} makes
this an exact equality of densities and of the ordinary and flagged
profiles. Its additional cost is at most $C_{\rm sec}p_\beta$, with
$C_{\rm sec}$ the fixed section and symmetry cost. Choose the probe
error to spend half the desired allowance $L^{-2/5}$ and the
nonprincipal error to spend the other half. The probe family and
frame are finite, so one entrance threshold meets all these demands.
This proves the small primary coefficient error for the prescribed
transport $x_i$, on the actual selected transition.

For completeness, its principal remainder has the required graded
bound as well. The local section for coordinate $i$ has zero moments
through degree $g_i-1$. Apply the short reference estimate with that
degree to its fixed stencil; it gives $C_iL^{-g_i}$ times the fine
mass. After $L$ is fixed, choose the comparison error so that its
reserved mass is at most $L^{-g_i}$ times that mass. Apply this to the
real and imaginary parts and multiply by $L^d$. The resulting
principal price is at most
\[
                     2(C_i+1)L^{d-g_i}m_i.
\]
Here the constant $C_i$ is purchased before the block factor using
the weighted response derivatives through order $g_i-1$.
The projection calculation already proved bounds the remaining cost
by fixed multiples of this price and $m_i$. Since $L^{d-g_i}\ge1$
and the normalized local-section mass $m_i$ is bounded by the fixed
frame credibility constant, choose the primary ledger constant $E_J$
larger than their finite sum over the frame. This gives the graded
remainder budget $E_JL^{d-g_i}$. It is the ordinary primary budget
used in the calibrated coefficient-field recursion, together with
the small coefficient error just proved.
